\chapter{Verbos modales básicos}
\section{Objetivos}
\begin{itemize}[leftmargin=*]
	\item Usar modales frecuentes en presente para expresar capacidad, permiso y deseo.
	\item Practicar infinitivos al final de la oración principal.
\end{itemize}

\section{Contenidos clave}
\begin{itemize}[leftmargin=*]
	\item Modales: \textit{kunnen, mogen, willen, moeten, zullen} (uso A1 de futuro/ intención con \textit{zullen}).
	\item Orden: [Sujeto] + [Modal conjugado] + ... + [Infinitivo final].
	\item Fórmulas de cortesía con \textit{mogen/kunnen}.
\end{itemize}

\section{Práctica sugerida}
\begin{itemize}[leftmargin=*]
	\item Completar oraciones con modal + infinitivo.
	\item Role-play de peticiones y permisos en tienda/cafetería.
	\item Transformar deseos en planes con \textit{zullen}.
\end{itemize}

\section{Microevaluación}
\begin{itemize}[leftmargin=*]
	\item 10 frases grabadas mezclando los cinco modales con verbos conocidos.
\end{itemize}
